\section{Bölüm sonu çözümlü problemler}

\begin{enumerate}
  \item \textbf{Tabakalı örnekleme.} Bir fakültede sınıf düzeylerine göre öğrenci sayıları 600, 500, 400 ve 300'dür. Orantılı tabakalı yöntemle $n=180$ öğrenci seçilecektir. Her sınıftan kaç öğrenci alınmalıdır?

  \textbf{Çözüm:} Toplam $N=1800$'dür. $n_h=n(N_h/N)$ kullanılır. Sırasıyla $180(600/1800)=60$, 50, 40 ve 30 öğrenci seçilir. Sayılar 180'e toplamlanır.

  \item \textbf{Ağırlıklı oran.} Örneklemde birinci sınıftan 80 öğrencinin 32'si, dördüncü sınıftan 20 öğrencinin 12'si danışma hizmetini kullanmıştır. Evrende iki sınıfın büyüklükleri eşitse ağırlıksız örneklem oranı ile sınıf ağırlıklı oranı bulun.

  \textbf{Çözüm:} Ağırlıksız oran $(32+12)/100=0{,}44$'tür. Sınıf oranları 0,40 ve 0,60'tır. Evren payları eşit olduğundan ağırlıklı oran $0{,}5(0{,}40)+0{,}5(0{,}60)=0{,}50$ olur. Orantısız seçim dikkate alınmadığında oran küçümsenmiştir.

  \item \textbf{Yanıtlamama.} 500 kişilik örneklemde 300 kişi ankete yanıt vermiş ve bunların yüzde 40'ı yüksek kaygı bildirmiştir. Yalnız bu bilgiyle evren oranı neden güvenle tahmin edilemez?

  \textbf{Çözüm:} Yanıt oranı yüzde 60'tır. Yanıtlamayanların kaygı düzeyi yanıtlayanlardan sistematik olarak farklıysa yüzde 40 yanlı olabilir. Hatırlatma dalgaları, bilinen yardımcı değişkenlerle ağırlıklandırma ve yanıtlayan--yanıtlamayan karşılaştırması gerekir; büyük örneklem yanıtlamama yanlılığını tek başına gidermez.
\end{enumerate}